% Part: computability
% Chapter: recursive-functions
% Section: trees

\documentclass[../../../include/open-logic-section]{subfiles}

\begin{document}

\olfileid{cmp}{rec}{tre}
\olsection{درخت‌ها}

گاهی سودمند است درخت‌ها را همچون اعداد طبیعی بازنمایی کنیم، همان‌گونه
که دنباله‌ها را با اعداد و ویژگی‌ها و عملیات روی آن‌ها را با روابط و
توابع بازگشتی اولیه روی کدهایشان بازنمایی می‌کنیم. برای این کار از
دنباله‌ها و کدهایشان استفاده خواهیم کرد. درخت می‌تواند یا تنها یک گره
(احتمالاً با برچسب) باشد، یا گرهی (احتمالاً با برچسب) که به چند
زیردرخت متصل است. آن گره 
\emph{ریشهٔ} درخت و زیردرخت‌هایی که به آن متصل‌اند 
\emph{زیردرخت‌های بی‌واسطهٔ} آن نامیده می‌شوند.

درخت‌ها را به‌طور بازگشتی به صورت دنبالهٔ
$\tuple{k,d_1,\dots,d_k}$ کدگذاری می‌کنیم که در آن~$k$ شمار
زیردرخت‌های بی‌واسطه و $d_1$، \dots، $d_k$ کدهای آن‌هاست. اگر گره‌ها
برچسب داشته باشند، می‌توان برچسب‌ها را پس از زیردرخت‌های بی‌واسطه
افزود. پس درختی که تنها از یک گره با برچسب~$l$ تشکیل شده باشد با
$\tuple{0,l}$ کدگذاری می‌شود؛ و درختی متشکل از ریشه‌ای با برچسب~$l_1$
که به دو گرهٔ منفرد با برچسب‌های~$l_2$ و~$l_3$ متصل است، با
$\tuple{2,\tuple{0,l_2},\tuple{0,l_3},l_1}$ کدگذاری می‌شود.

\begin{prop}
  \ollabel{prop:subtreeseq}
  تابع~$\fn{SubtreeSeq}(t)$ که کد دنباله‌ای را بازمی‌گرداند که عناصرش
  کدهای همهٔ زیردرخت‌های درخت با کد~$t$ هستند، بازگشتی اولیه است.
\end{prop}

\begin{proof}
  نخست توجه کنید که
  $\fn{ISubtrees}(t)=\fn{subseq}(t,1,(t)_0)$ بازگشتی اولیه است و
  کدهای زیردرخت‌های بی‌واسطهٔ درخت~$t$ را بازمی‌گرداند. اکنون می‌توانیم
  تابع کمکی~$\fn{hSubtreeSeq}(t,n)$ را تعریف کنیم که دنبالهٔ همهٔ
  زیردرخت‌هایی را محاسبه می‌کند که فاصله‌شان از ریشه حداکثر~$n$ گره
  است. دنبالهٔ زیردرخت‌های~$t$ در فاصلهٔ حداکثر صفر از ریشه---به بیان
  دیگر، دنباله‌ای که از خود ریشهٔ~$t$ آغاز می‌شود---تنها از~$t$ تشکیل
  شده است. برای به‌دست‌آوردن دنبالهٔ زیردرخت‌ها تا سطح~$n+1$، دنبالهٔ
  تا سطح~$n$ را با دنبالهٔ همهٔ زیردرخت‌های بی‌واسطهٔ اعضای آن الحاق
  می‌کنیم. برای فهرست‌کردن همهٔ این زیردرخت‌ها، توجه کنید اگر~$f(x)$
  تابعی بازگشتی اولیه باشد که کد دنباله‌ها را بازمی‌گرداند، تابعی که
  مقادیر~$f$ را برای نخستین~$k$ عنصر~$s$ به هم الحاق می‌کند نیز
  بازگشتی اولیه است:
    \begin{align*}
      g_f(s, 0) & = \emptyseq\\
      g_f(s, k+1) & = g_f(s, k) \concat f((s)_k).
    \end{align*}
    برای نمونه، اگر~$s$ دنباله‌ای از درخت‌ها باشد،
    $h(s)=g_{\fn{ISubtrees}}(s,\len{s})$ دنبالهٔ زیردرخت‌های بی‌واسطهٔ
    عناصر~$s$ را به دست می‌دهد. با استفاده از آن،
    $\fn{hSubtreeSeq}$ را چنین تعریف می‌کنیم:
    \begin{align*}
      \fn{hSubtreeSeq}(t, 0) & = \tuple{t} \\
      \fn{hSubtreeSeq}(t, n+1) & = \fn{hSubtreeSeq}(t, n) \concat
      h(\fn{hSubtreeSeq}(t, n)).
    \end{align*}
    بیشینهٔ سطح زیردرخت‌ها در درختی با کد~$t$، یعنی بیشینهٔ فاصلهٔ
    ریشه تا یک گرهٔ برگ، با خود کد~$t$ کراندار است. پس دنباله‌ای از
    کدهای همهٔ زیردرخت‌های درخت با کد~$t$ را
    $\fn{hSubtreeSeq}(t,t)$ به دست می‌دهد.
\end{proof}

\begin{prob}
  تعریف~$\fn{hSubtreeSeq}$ در اثبات
  \olref[cmp][rec][tre]{prop:subtreeseq} به‌طور کلی شامل تکرار است.
  تعریفی جایگزین ارائه کنید که تضمین کند کد هر زیردرخت تنها یک بار در
  فهرست حاصل رخ می‌دهد.
\end{prob}

\end{document}

